\subsection{Treewidth}
\label{sec:tw-transfer}

We prove UGC-hardness below $4/3$ by transferring a balanced,
unweighted tripartite Vertex Cover gap on graphs of bounded degree
to treewidth. We derive this gap from
Theorem~\ref{thm:bounded-degree-kr-source} using a multipartite product.
B\'erczi, Chandrasekaran, Kir\'aly, Lee, and Xu prove a related weighted
$k$-partite Vertex Cover gap under UGC~\cite[Sections~3.5--3.6]{BCKLX2017}.
The same transfer gives the NP-hardness bound in
Section~\ref{subsec:np-treewidth}.

The transfer uses Bonnet's completion of three parts to cliques
and his central-bag lower bound~\cite[Lemma~6]{BonnetTreewidth}.
We take disjoint source copies \emph{before} completing the three global parts.
This lets each of three paths attached to a central bag process
the copies one at a time: the additional bag size is bounded by one part of one copy,
while the cover contribution grows with the number of copies.
Completing the global parts creates a dense intermediate graph.
Corollary~\ref{cor:tw-bounded-degree} obtains bounded-degree output
by replacing these three cliques with bounded-degree expanders while
retaining all source-copy edges.

\subsubsection{A transfer from vertex cover to treewidth}

Let $G$ be a nonempty tripartite graph with specified independent
parts $V(G)=A\mathbin{\dot\cup}B\mathbin{\dot\cup}C$, and put
$n_{\max}=\max\{|A|,|B|,|C|\}$. Thus $G$ has no edge within a part.
For an integer $R\ge1$, take copies $G_1,\ldots,G_R$ on
$V(G)\times\{1\},\ldots,V(G)\times\{R\}$, writing $v_j=(v,j)$.
For $X\in\{A,B,C\}$ and $j\in[R]$, let $X_j=X\times\{j\}$ and
$X^*=X\times[R]$. Define $K_R(G)$ on $V(G)\times[R]$ by
writing $\binom{X^*}{2}$ for the set of unordered pairs of distinct
vertices in $X^*$:
\[
 E(K_R(G))=
 \bigcup_{X\in\{A,B,C\}}\binom{X^*}{2}
 \;\cup\;\bigcup_{j=1}^R\{\{u_j,v_j\}:uv\in E(G)\}.
\]
These are all the edges: each global part is a
clique, and vertices in different global parts are adjacent if and only if
they belong to the same copy and their source vertices are
adjacent in $G$.

Its vertex partition into three cliques makes $K_R(G)$
\emph{co-tripartite}.

\begin{lemma}[A bag covering the source edges]\label{lem:tw-central-bag}
Every tree decomposition of $K_R(G)$ has a bag containing a vertex
cover of $G_1\mathbin{\dot\cup}\cdots\mathbin{\dot\cup}G_R$.
In particular, that bag has size at least $R\vc(G)$.
\end{lemma}

\begin{proof}
We use Bonnet's central-bag argument
\cite[proof of Lemma~6]{BonnetTreewidth}.
Fix a tree decomposition $(T,(B_t)_{t\in V(T)})$ of $K_R(G)$.
For each nonempty clique $X^*$, the subtrees $T_x$, $x\in X^*$,
intersect pairwise. By Lemma~\ref{lem:subtree-helly},
they have a common node $u_X$, so $X^*\subseteq B_{u_X}$
\cite[Lemma~4]{BonnetTreewidth}. Choose $u_X\in V(T)$ arbitrarily if
$X^*=\varnothing$. Lemma~\ref{lem:three-tree-paths} gives a unique
common node $u$ of the three paths joining pairs of
$u_A,u_B,u_C$.
For an edge $xy\in E(G_j)$, with $j\in[R]$, let
$x\in X^*$ and $y\in Y^*$ for distinct $X,Y\in\{A,B,C\}$.
The subtrees $T_x,T_y$ intersect. Their union is
connected and contains $u_X,u_Y$, so it contains the
$u_X$--$u_Y$ path and hence $u$. Thus $u\in T_x\cup T_y$,
which means that $B_u$ contains $x$ or $y$.
Therefore $B_u$ contains an endpoint of every source edge,
so it contains a vertex cover of the disjoint source copies
and has size at least $R\vc(G)$.
\end{proof}

\begin{lemma}[A decomposition from a source cover]\label{lem:tw-cover-decomposition}
For every vertex cover $S$ of $G$, the graph $K_R(G)$ has a tree
decomposition whose bags have size at most
$R\max\{|S|,n_{\max}\}+n_{\max}$.
\end{lemma}

\begin{proof}
We adapt the three-path decomposition of~\cite[Section~3.3]{BonnetTreewidth}
by processing copies separately, so the temporary bag increase is at
most $n_{\max}$, independently of $R$.
Put $I=V(G)\setminus S$.
For each $j\in[R]$, put $S_j=S\times\{j\}$ and $I_j=I\times\{j\}$.
Start with the central bag
$S^*=\bigcup_{j=1}^R S_j$, where $|S^*|=R|S|$. Attach three paths to this
node, one for each $X\in\{A,B,C\}$. The path to $X$ ends with the bag $X^*$.
On this path process the copies in order $j=1,\ldots,R$. At step $j$,
first add all vertices of $I_j\cap X_j$, and then remove all vertices
of $S_j\setminus X_j$.

After processing the first $j$ copies, the bag is
\[
 B^X_j=\bigcup_{i=1}^j X_i\;\cup\!\bigcup_{i=j+1}^R S_i
 \qquad(0\le j\le R).
\]
Here $B^X_0=S^*$ and $B^X_R=X^*$. Between $B^X_{j-1}$ and
$B^X_j$ insert the bag $B^X_{j-1}\cup(I_j\cap X_j)$.
The bag sizes satisfy
\[
\begin{aligned}
 |B^X_j|&=j|X|+(R-j)|S|\le R\max\{|S|,|X|\},\\
 |B^X_{j-1}\cup(I_j\cap X_j)|
 &=|B^X_{j-1}|+|I\cap X|
 \le R\max\{|S|,n_{\max}\}+n_{\max}.
\end{aligned}
\]
We verify the decomposition axioms. Every vertex is present in its
part's final bag. A vertex in $S_j\cap X_j$ occurs at the center and
throughout the $X$-path, and on each other path it is removed once.
A vertex in $I_j\cap X_j$ occurs only on the $X$-path, from its
addition to the end. Thus each $T_v$ is nonempty and connected.
Every edge within a global clique is covered at
the corresponding final bag. Now consider an edge between two global
parts. It comes from a single source copy, and its endpoints cannot
both lie in that copy's independent set. If both endpoints lie in
$S^*$, the center bag covers the edge. Otherwise, one endpoint lies
in $I_j\cap X_j$ and the other in $S_j\setminus X_j$ for some $j$
and part $X$. When the first endpoint is added on the $X$-path, the
second is still in the bag, so the edge is covered there. All edges
are covered, and the bag-size bound follows from the count above.
\end{proof}

\begin{lemma}[Replicated co-tripartite completion]\label{lem:tw-replication}
Let $G$ be a nonempty tripartite graph with specified independent
parts $A,B,C$, put $n_{\max}=\max\{|A|,|B|,|C|\}$, and let
$R\ge1$ be an integer. Then
\[
 R\max\{\vc(G),n_{\max}\}-1
 \leq\operatorname{tw}(K_R(G))\leq
 R\max\{\vc(G),n_{\max}\}+n_{\max}-1.
\]
\end{lemma}


\begin{proof}
Lemma~\ref{lem:tw-central-bag} gives a bag of size at least $R\vc(G)$,
and Lemma~\ref{lem:subtree-helly} applied to a largest global clique
gives one of size at least $Rn_{\max}$.
Thus every decomposition has a bag of size at least
$R\max\{\vc(G),n_{\max}\}$. For the upper bound, apply
Lemma~\ref{lem:tw-cover-decomposition} to a minimum cover of $G$:
its largest bag has size at most
$R\max\{\vc(G),n_{\max}\}+n_{\max}$. Subtracting one from both
bag-size bounds proves the claim.
\end{proof}

For fixed $R$, the construction has $R|V(G)|$ vertices and is
polynomial in the source size.

\subsubsection{The approximation gap}
\label{subsec:ugc-tw-gap}

For the UGC reduction we need cover thresholds $(1/2+\zeta)N$ and
$(2/3-\zeta)N$ on balanced tripartite graphs.
We derive the balanced, unweighted form needed here from the
bounded-degree independent-set promise using the product below. Its degree
identity will also be used in
Section~\ref{sec:bounded-degree-hardness}.

\begin{lemma}[Multipartite product]\label{lem:multipartite-product}
For a nonempty simple graph $G$ on $n$ vertices and a positive integer
$k$, let $P_k(G)$ have vertex set $V(G)\times[k]$, and write
$v_i=(v,i)$. For $u,v\in V(G)$ and $i,j\in[k]$, define its edges by
\[
 u_iv_j\in E(P_k(G))\quad\Longleftrightarrow\quad
 i\ne j\text{ and }uv\in E(G).
\]
Its layers $V(G)\times\{i\}$, $i\in[k]$, are independent sets of size $n$. Moreover,
$k\alphaI(G)\le\alphaI(P_k(G))\le n+(k-1)\alphaI(G)$ and
$\Delta(P_k(G))=(k-1)\Delta(G)$.
\end{lemma}

\begin{proof}
All $k$ copies of an independent set of $G$ form an independent
set of $P_k(G)$, proving the first bound. For the upper bound, let
$I$ be any independent set in the product. Put
$m_v=|\{i\in[k]:v_i\in I\}|$ for $v\in V(G)$ and
$D=\{v\in V(G):m_v\ge2\}$. If distinct $u,v\in D$
were adjacent in $G$, choose a copy $u_i\in I$. The set $I$ contains
at least two copies of $v$, so one has a layer index $j\ne i$. Then
$u_iv_j$ would be an edge inside $I$. Thus $D$ is independent, so
\begin{equation}\label{eq:multipartite-repeated-support}
 |D|\leq\alphaI(G).
\end{equation}
Since $m_v\leq k$ for $v\in D$ and $m_v\leq1$ otherwise,
\[
 |I|=\sum_{v\in V(G)}m_v\leq n+(k-1)|D|
 \overset{\eqref{eq:multipartite-repeated-support}}{\leq}
 n+(k-1)\alphaI(G).
\]
For every $v\in V(G)$ and $i\in[k]$,
$\deg_{P_k(G)}(v_i)=(k-1)\deg_G(v)$,
which proves the maximum-degree identity.
\end{proof}

\begin{theorem}[Balanced bounded-degree multipartite cover gap]
\label{thm:ugc-tripartite-cover}
For every fixed integer $k\geq3$ and fixed
$0<\zeta<1/12$, there is a constant $D_{k,\zeta}$ such that the
following promise problem is Unique Games-hard. The input is a
simple unweighted $k$-partite graph $G$ on $N\geq k$ vertices, of maximum
degree at most $D_{k,\zeta}$, with a specified partition into $k$
independent parts each of size $N/k$. Distinguish the following cases:
\textbf{YES:} $\vc(G)\leq(1/2+\zeta)N$;
\textbf{NO:} $\vc(G)\geq(1-1/k-\zeta)N$.
The reduction is deterministic and polynomial time.
\end{theorem}

The promise follows from Theorem~\ref{thm:bounded-degree-kr-source} and
Lemma~\ref{lem:multipartite-product}, with
$D_{k,\zeta}=(k-1)d_{\zeta/2}$. The independent-set source is already
of bounded degree by the spectral sparsification in
Appendix~\ref{app:bounded-degree-ugc-source}.
Appendix~\ref{app:multipartite-normalization} records the parameter
substitution. Taking $k=3$ and applying the
replication lemma proves the treewidth gap.



\begin{proof}[Dense treewidth gap: proof sketch]
Apply Theorem~\ref{thm:ugc-tripartite-cover} with $k=3$ and then
Lemma~\ref{lem:tw-replication}. For the $N$-vertex source, the
YES and NO width bounds are respectively
\[
 R(1/2+\zeta)N+N/3-1,\qquad R(2/3-\zeta)N-1.
\]
As the source error decreases and the fixed replication count $R$
increases, their leading coefficient ratio approaches $4/3$.
Thus every fixed $r<4/3$ is excluded. The same comparison applies
to maximum bag size by adding one to both bounds.
Appendix~\ref{app:dense-ugc-tw} supplies the choices of $\zeta,R$
and checks the additive constants and decision thresholds.
Writing $\Lambda$ for the Unique Games instance, the construction is
\[
 \Lambda\xrightarrow{\text{tripartite source}}G
 \xrightarrow{\text{replicate and complete}}K_R(G).
\]
\end{proof}
